%==============================================================================
\section{Khái niệm nền --- Hồi quy trong học máy}
\label{sec:khai-niem-hoi-quy}
%==============================================================================

\begin{dinhnghia}[title={Phân tích hồi quy (Regression Analysis)}]
Kỹ thuật thống kê dự đoán \textbf{giá trị số liên tục} dựa trên quan hệ giữa biến độc lập và biến phụ thuộc.
Mục tiêu: vẽ đường/đường cong khớp dữ liệu tốt nhất và ước lượng ảnh hưởng của biến này lên biến kia.
\end{dinhnghia}

\begin{ynghia}[title={Học có giám sát --- hồi quy}]
Đầu ra là \textbf{số liên tục}; mô hình học từ feature (độc lập) và target (phụ thuộc) để dự đoán trên dữ liệu mới.
\end{ynghia}

\begin{phuongphap}[title={Thuật ngữ thường gặp}]
\begin{itemize}
  \item \textbf{Biến độc lập} (predictor, feature).
  \item \textbf{Biến phụ thuộc} (target, response).
  \item \textbf{Đường hồi quy}: đường/đường cong khớp dữ liệu.
  \item \textbf{Overfitting / Underfitting}: variance cao / bias cao.
  \item \textbf{Outlier}, \textbf{đa cộng tuyến} (multicollinearity).
\end{itemize}
\end{phuongphap}

\begin{phuongphap}[title={Các chương trong phần này}]
Hồi quy tuyến tính (tổng quan), hồi quy tuyến tính đơn, hồi quy tuyến tính bội, hồi quy đa thức.
\end{phuongphap}
